% Generated draft: requires mathematical and source review.
\documentclass[11pt]{article}
\usepackage[margin=25mm]{geometry}
\usepackage{fontspec}
\setmainfont{Latin Modern Roman}
\usepackage{amsmath,amssymb,mathtools,unicode-math}
\setmathfont{Latin Modern Math}
\usepackage{enumitem,needspace,xurl,hyperref}
\setlength{\parindent}{0pt}
\setlength{\parskip}{5pt}
\setlength{\emergencystretch}{3em}
\newcommand{\N}{\mathbb{N}}
\newcommand{\Z}{\mathbb{Z}}
\newcommand{\Q}{\mathbb{Q}}
\newcommand{\R}{\mathbb{R}}
\newcommand{\C}{\mathbb{C}}
\newcommand{\F}{\mathbb{F}}
\newcommand{\A}{\mathbb{A}}
\newcommand{\PP}{\mathbb{P}}
\newcommand{\E}{\mathbb{E}}
\newcommand{\eps}{\varepsilon}
\newcommand{\dd}{\,\mathrm{d}}
\begin{document}
\section*{UnsolvedMath ID: 800655}
\textbf{Source problem code:} catalogue entry
\subsection{Definitions and mathematical statement}
For families \(\mathcal A\) and \(\mathcal B\) of covers of a space \(X\), define
\[
\binom{\mathcal A}{\mathcal B}
\quad\Longleftrightarrow\quad
(\forall\mathcal U\in\mathcal A)(\exists\mathcal V\subseteq\mathcal U)\;\mathcal V\in\mathcal B.
\]
In standard selection-principle terminology, an open \(\omega\)-cover of \(X\) is a family \(\mathcal U\) of open subsets of \(X\) such that \(X\notin\mathcal U\) and every finite subset of \(X\) is contained in some member of \(\mathcal U\). An open \(\gamma\)-cover is an infinite open cover \(\mathcal U\) such that each point of \(X\) belongs to all but finitely many members of \(\mathcal U\). The symbols \(\Omega\) and \(\Gamma\) conventionally denote the corresponding families. The supplied record does not specify countability restrictions or the full cover conventions, so their exact source-specific meanings require confirmation.

Let \(\mathscr P\) stand for the intended collection of ultrafilters. The supplied record does not specify their underlying set, whether they must be free, or the definition of the associated cover family \(\mathcal O_p\). With these unresolved source parameters retained, the question is whether every space \(X\) in the intended class satisfies
\[
\left[(\forall p\in\mathscr P)\;\binom{\Omega(X)}{\mathcal O_p(X)}\right]
\quad\Longrightarrow\quad
\binom{\Omega(X)}{\Gamma(X)}.
\]
Equivalently, the hypothesis says
\[
(\forall p\in\mathscr P)(\forall\mathcal U\in\Omega(X))
(\exists\mathcal V\subseteq\mathcal U)\;\mathcal V\in\mathcal O_p(X),
\]
and the requested conclusion is
\[
(\forall\mathcal U\in\Omega(X))
(\exists\mathcal V\subseteq\mathcal U)\;\mathcal V\in\Gamma(X).
\]
Thus the selected subcover in the hypothesis may depend on both \(p\) and \(\mathcal U\). The intended class of spaces and the missing ultrafilter-cover definitions must be recovered before this schematic formulation becomes a fully specified mathematical statement.
\subsection{Short English statement}
Must a space having an \(\mathcal O_p\)-subcover of every \(\omega\)-cover, separately for every intended ultrafilter \(p\), have a \(\gamma\)-subcover of every \(\omega\)-cover? The space class, ultrafilter domain, and precise meaning of \(\mathcal O_p\) are unspecified in the supplied record.
\subsection{Sources}
\begin{itemize}
\item Supplied local UnsolvedMath registry record: permanent ID \(800655\), source code \texttt{catalogue entry}, title ``Selection Principles Milestone Problems --- Problem 28 ([GTpseq, Problem 3.14(2)]).'' The registry lists Tsaban as proposer, gives the year \(2003\), and labels the problem open; this status is supplied metadata, not an independently verified literature assessment.
\item Supplied source: Tsaban, \emph{Selection Principles in Mathematics: A Milestone of Open Problems} (2003), Problem 28, \url{https://arxiv.org/abs/math/0312182}. The record reports access on 2026-10-16 and extraction by \texttt{arxiv\_latex}; these are supplied provenance statements, not a claim of independent source consultation.
\item Supplied attribution/context: [GTpseq, Problem 3.14(2)]. The record does not supply the bibliographic expansion of \texttt{GTpseq}.
\end{itemize}
\end{document}
